\documentclass[11pt]{article}
\usepackage[T1]{fontenc}
\usepackage{lmodern,amsmath,amssymb,amsthm,geometry,microtype}
\usepackage[hidelinks]{hyperref}
\geometry{margin=1in}
\newtheorem{theorem}{Theorem}
\newtheorem{lemma}{Lemma}
\newcommand{\R}{\mathbb R}
\newcommand{\T}{\mathbb T}
\newcommand{\Id}{\mathrm{Id}}
\title{A torus-collar proof of the real-spectrum compact-support classification for Nijenhuis operators}
\author{Alec Kriebel\\Independent researcher\\\small ORCID: \href{https://orcid.org/0009-0001-9320-500X}{0009-0001-9320-500X}}
\date{15 October 2026, version 1.0}
\begin{document}
\maketitle
\begin{abstract}
We give a complete smooth compact-support proof of the classification asked
for in the historical integrable-systems problem source: a real Jordan type
occurs on a nonempty open subset of $\R^n$ for a compactly supported smooth
Nijenhuis operator exactly when its spectrum is real. The construction
realizes arbitrary specified numerical real Jordan data in a commuting local
frame. On a same-dimensional torus collar, the torsion equation reduces to
$(A-\lambda I)A'=0$; an ordered finite path removes every Jordan block and
reaches an exact zero band. The nonreal obstruction and earlier semisimple
and two-dimensional examples are prior results. We include the precise
zero-point rigidity deduction needed for necessity. The historical priority
of the full result and of this construction remains unestablished.
Extensive AI tools assisted solving, verification and writing. This preprint
is unrefereed and has no external human peer review or formal
proof-assistant certification.
\end{abstract}

\section{Statement, historical target and attribution}
For a smooth real field $L$ of tangent-space endomorphisms, use the convention
\[
 N_L(X,Y)=L^2[X,Y]+[LX,LY]-L[LX,Y]-L[X,LY].
\]
An operator is Nijenhuis when $N_L=0$. Jordan type here means the block
sizes and grouping by eigenvalue; our existence statement also fixes any
desired numerical real eigenvalues.

\begin{theorem}\label{thm:main}
For every $n\geq1$ and every prescribed real Jordan matrix $J$ with entirely
real spectrum, there exists a globally $C^\infty$, compactly supported real
Nijenhuis operator on $\R^n$ whose matrix is $J$ in a smooth commuting local
frame on a nonempty open domain. Conversely, every eigenvalue at every
point of a compactly supported smooth real Nijenhuis operator on $\R^n$ is
real.
\end{theorem}

The type, rank and eigenvalues may change outside the realization domain.
The theorem prescribes neither the matrix in the ambient Cartesian frame
nor global regularity, and is a smooth, not a real-analytic, statement.
The zero matrix is permitted.

The target is Problem 5.11 on p.~29 of the 14 January 2020 author-v2
of Bolsinov--Matveev--Miranda--Tabachnikov~\cite{problems}; its 10 April
2018 author-v1 counterpart is Problem 4.10 on p.~37. The original journal
work appeared in 2018. A corpus label ``5.9'' is not the locator in either
inspected version. The full wording asks for necessary and sufficient
conditions for compact support with a prescribed unchanged algebraic type
on a domain. Theorem~\ref{thm:main} answers that printed question, without
asserting that it remained open worldwide before this note.

The source already credits the nonreal obstruction, real semisimple
realizations and a two-dimensional Jordan example. Bolsinov--Konyaev--Matveev
(BKM)~\cite{bkm}, author-v2 Theorem 4.5 and the following p.~28 remark,
also give arbitrary algebraic types at local points. Those facts are prior.
The proof below supplies an explicit fixed-type open domain together with
a smooth compact-support transition. The bounded literature audit in the
supplement certifies neither firstness of the full theorem nor priority of
the torus-collar mechanism.

\section{The one-variable collar equation}
Let $m=n-1\geq1$ and work on $\T^m\times(-1,1)$. Write
$e_1,\ldots,e_m$ for the invariant angular fields and $v=\partial_t$.
They form a global commuting frame. Angular coordinate functions are only
locally real-valued; the invariant fields and their dual forms are global.
For smooth matrix, vector and scalar functions of $t$, set
\[
 Le_i=\sum_{k=1}^m A_{ki}(t)e_k,\qquad
 Lv=\sum_{k=1}^m b_k(t)e_k+\lambda(t)v.
\]
All angular-pair torsions vanish. For a mixed pair the brackets are
\[
 [Le_i,Lv]=-\lambda A'e_i,\quad [Le_i,v]=-A'e_i,
 \quad [e_i,Lv]=[e_i,v]=0.
\]
Consequently
\begin{equation}\label{eq:collar}
 N_L(e_i,v)=(A-\lambda I)A'e_i,\qquad
 N_L=0\ \Longleftrightarrow\ (A-\lambda I)A'=0.
\end{equation}
Tensoriality and skew symmetry account for every component. The product
order is essential. There is no condition on $b'$, $\lambda'$, rank,
diagonalizability or invertibility.

\begin{lemma}\label{lem:path}
For every real-spectrum Jordan matrix $J$ of size $n\geq2$, after a constant
permutation of whole blocks there is a smooth path
$L(s)=\left(\begin{smallmatrix}A(s)&b(s)\\0&\lambda(s)\end{smallmatrix}\right)$,
$0\leq s\leq1$, from $J$ to zero, constant near both endpoints, satisfying
$(A-\lambda I)\,dA/ds=0$.
\end{lemma}
\begin{proof}
Use upper-superdiagonal Jordan convention and designate the last vector of
one block as transverse. The remaining hyperplane is invariant. Thus
$J=\left(\begin{smallmatrix}A_0&b_0\\0&\lambda_0\end{smallmatrix}\right)$,
where $A_0$ contains the other blocks and the distinguished block shortened
by one; a size-one distinguished block disappears. The vector $b_0$ retains
the removed terminal arrow, so the initial full matrix is exactly $J$.

First hold $A$ and $\lambda$ fixed and move $b_0$ to zero. Whenever $A$ is
fixed, both $b$ and $\lambda$ can move freely in~\eqref{eq:collar}.
Process the blocks of $A_0$ one at a time, leaving processed blocks zero and
unprocessed blocks fixed. For a current block $\mu I_r+K$, move $\lambda$
to the real number $\mu$ while holding $A$ fixed. Write
\[
 K=\sum_{j=1}^{r-1}a_jE_{j,j+1},\qquad
 KK'=\sum_{j=1}^{r-2}a_j a'_{j+1}E_{j,j+2}.
\]
Initially every $a_j=1$. With $\lambda=\mu$, move $a_1$ to zero, then
$a_2$ to zero, and so forth. Only one arrow changes at a time. Whenever
$a'_{j+1}\ne0$, its predecessor $a_j$ is already identically zero; the
first arrow has no predecessor. Hence $KK'=0$. Other blocks have zero
derivative and contribute nothing to the ordered product.

The active block is now scalar. Move its scalar value $c(s)$ from $\mu$
to zero and set $\lambda(s)=c(s)$ on this stage. The derivative is supported
in that block, where $A-\lambda I$ vanishes. Repeat for every angular block.
For singleton blocks there are no arrows; zero eigenvalues allow constant
scalar stages. At the end $A=b=\lambda=0$.

There are finitely many stages. On separate successive intervals use a
$C^\infty$ switch constant near both endpoints. For example, with
$H(u)=e^{-1/u}$ for $u>0$ and $H(u)=0$ otherwise, the function
$S(u)=H(u)/(H(u)+H(1-u))$ is smooth, zero for $u\leq0$ and one for
$u\geq1$; apply $S$ to an affine parameter whose endpoints extend beyond
$[0,1]$. Each stage then has genuine open endpoint plateaus. Values agree
at joins, so concatenation is smooth. Reparametrization only multiplies
$A'$ by a scalar, preserving~\eqref{eq:collar}. No division by an eigenvalue
or eigenvalue difference is used.
\end{proof}

Choose a smooth even $q:(-1,1)\to[0,1]$ equal to zero for $|t|\leq1/4$
and one for $|t|\geq3/4$; explicitly take
$q(t)=S(2(t^2-1/16))$. Substitute $s=q(t)$ in Lemma~\ref{lem:path}.
The chain rule gives
\[
 (A(q)-\lambda(q)I)\frac{d}{dt}A(q)
 =q'(t)(A(q)-\lambda(q)I)\frac{d}{ds}A(q)=0.
\]
The tensor is $J$ on the central band and exactly zero on both end bands.
There is no absolute-value corner and no arbitrary scalar cutoff of $J$.
A constant frame permutation restores any initially specified block order;
constant linear combinations of commuting fields still commute.

\section{A same-dimensional torus collar and extension}
For every $d\geq1$, the torus $\T^d$ embeds as a two-sided smooth
hypersurface in $\R^{d+1}$. Here is the needed all-dimensional induction.
Start with a circle in $\R^2$. If $f:\T^{d-1}\hookrightarrow\R^d$ is an
oriented embedded hypersurface, it has a global unit normal $\nu$.
Inside $\R^{d+1}$ its rank-two normal bundle has the explicit frame
$(\nu,0),(0,1)$. By compactness and the tubular-neighborhood theorem, a
sufficiently small normal disk bundle is embedded. Its boundary is
$\T^{d-1}\times S^1=\T^d$, with parametrization
\[
 F(p,\phi)=\big(f(p)+\varepsilon\cos\phi\,\nu(p),
                         \varepsilon\sin\phi\big).
\]
The tubular coordinates recover $p,\cos\phi,\sin\phi$, proving injection
and immersion. Compactness makes this an embedding. A global unit normal
is $(\cos\phi\,\nu(p),\sin\phi)$: orthogonality follows from
$\nu\cdot df=\nu\cdot d\nu=0$. This verifies two-sidedness for the next
step and completes the induction.

For $n\geq2$ a thin normal collar of the embedded $\T^{n-1}\subset\R^n$
is diffeomorphic, after rescaling, to $\T^{n-1}\times(-1,1)$. Push the
preceding tensor through this diffeomorphism. Its frame still commutes,
and torsion vanishes by naturality. Its support lies in the image of the
compact set $\T^{n-1}\times[-3/4,3/4]$, strictly inside the open collar.
Extend it by zero to $\R^n$. Each collar-boundary or exterior point has
a neighborhood outside that compact support. Thus the extension is
globally smooth, and its torsion is zero by locality. No exterior angular
frame is required. The central band is nonempty and open; a connected
coordinate box inside it is a realization domain.

For $n=1$, a compactly supported smooth scalar function equal to the desired
eigenvalue on an interval, multiplied by $\Id$, suffices. The alternating
vector-valued two-form $N_L$ vanishes in dimension one. The zero matrix is
also realized by the zero tensor. This proves the existence direction in
every stated boundary case.

\section{The credited nonreal obstruction}
BKM's author-v2 Theorem 6.1~\cite{bkm}, pp.~41--43, states nonreal spectral
persistence with constant algebraic multiplicity on a closed connected
smooth manifold. It has no analyticity or regularity hypothesis. The
following weaker consequence is all that is needed here. We give its
maximum-principle justification using their local splitting and canonical
complex/holomorphic results (Theorems 3.1 and 3.3).

\begin{lemma}\label{lem:zero}
A smooth real Nijenhuis operator on a closed connected manifold which
vanishes at one point has entirely real spectrum everywhere.
\end{lemma}
\begin{proof}
Suppose otherwise. Let $\mu(x)$ be the largest positive imaginary part of
an eigenvalue, with value zero for an all-real spectrum. Roots depend
continuously on polynomial coefficients, so $\mu$ is continuous. On the
compact manifold it has a positive maximum $\mu_0$. Among roots on the
line $\operatorname{Im}z=\mu_0$ at maximum points choose the largest
algebraic multiplicity $m_0$. The locus $E$ having such a root of
multiplicity $m_0$ is nonempty and closed: along a convergent sequence
choose a convergent subsequence of those roots; vanishing polynomial
derivatives shows multiplicity at least $m_0$, and maximality gives equality.

At $p\in E$, isolate the root $\lambda_0$ and its conjugate by the local
splitting theorem. On this real block the canonical complex structure
makes its complex characteristic factor
$P(z)=z^{m_0}+a_1z^{m_0-1}+\cdots+a_{m_0}$ holomorphic in the block's
complex coordinates and independent of complementary coordinates. At $p$,
$P(z)=(z-\lambda_0)^{m_0}$. All nearby roots have imaginary part at most
$\mu_0$. Therefore $-\operatorname{Im}a_1\leq m_0\mu_0$ with equality
at $p$. The holomorphic maximum principle makes $a_1$ constant locally.
Every root consequently has imaginary part $\mu_0$. The shifted polynomial
$P(w+i\mu_0)$ has only real roots. Its coefficients are real-valued
holomorphic functions and hence constant. Thus $P$ remains
$(z-\lambda_0)^{m_0}$ locally. The complementary block can be empty;
otherwise independence of its coordinates gives the same conclusion on
the product neighborhood. Hence $E$ is open. Connectedness gives $E=M$,
contradicting the zero point.
\end{proof}

For compactly supported $L$ on $\R^n$, identify $\R^n$ with a standard
chart of $S^n$ and extend by zero near infinity. Compact support makes this
a smooth Nijenhuis field on closed connected $S^n$, zero on a neighborhood
of infinity. Lemma~\ref{lem:zero} excludes nonreal spectrum, completing
Theorem~\ref{thm:main}.

The attribution is substantive: both the obstruction and its local
splitting/holomorphic mechanism are prior. Our argument avoids relying on
the final polynomial-reduction sentence of the author-v2 proof, whose
literal count of removed nonreal pairs is imprecise. For example, rotation
blocks with spectra $\{\pm i,\pm i/2\}$ all acquire real eigenvalues under
$z\mapsto z^2+1$, although two original pairs were nonreal. This is not a
counterexample to the theorem; it explains why the exact application is
proved by the preceding maximum argument. We do not certify every step of
the stronger persistence proof or identify the unread Advances in
Mathematics publisher body with author-v2 bytes.

\section{Relation to prior constructions and verification}
The arbitrary-type point construction of BKM Theorem 4.5 and its remark
is generically semisimple in independent characteristic-coefficient
coordinates. Repeated-root conditions there do not supply a fixed
nonsemisimple type on an open domain. In Geometry III~\cite{bkm3}, the
actual journal regular-layer theorem is Theorem 1.10 (author-v1 Theorem
1.4). It realizes regular layers near a nonzero nilpotent block, rather
than arbitrary repeated blocks with a compact zero transition. The journal
and arXiv numberings must be distinguished.

Geometry IV~\cite{bkm4} does allow singular or zero output symmetries of a
regular generator. Its functional coefficients nevertheless obey
compatibility equations. In its size-three nilpotent formula an unwrapped
angular coordinate multiplies a transverse derivative, so naive
periodization does not provide the collar construction. This observation
does not exclude a different use of that framework. Local constant-type
or spectral splitting results and fixed functional calculus likewise do
not by themselves supply compact support: a product block sum retains
unbounded cylinders of support, and fixed calculus preserves
diagonalizability of a semisimple input. Standard tubular-neighborhood and
smooth-switch arguments are not claimed as new.

The accompanying portable standard-library controls check the full collar
identity from coordinate first jets, finite ordered Jordan schedules with
exact polynomial arithmetic, and deliberately invalid arrow order,
cutoff and unmatched-scalar variants. They corroborate indexing; the
all-dimensional proof is above. These are newly distributed controls, not
a reconstruction or certification of the submitted historical omitted
programs or their counts. The source ledger records actual prior reading
scopes and unknowns, including unread 2025 Sbornik and 2026 cohomology
bodies and older cited works. No exhaustive historical inference is made.

\paragraph{Disclosure and license.}
AI tools were used extensively in solving, source comparison, adversarial
verification and writing. This note is unrefereed; automated and AI reviews
are not external human peer review. No journal acceptance or formal
proof-assistant certification is claimed. Full-result and construction-method
priority remain unknown. Copyright Alec Kriebel, 2026; distributed under
Creative Commons Attribution 4.0 International (CC BY 4.0).

\begin{thebibliography}{9}
\bibitem{problems} A. Bolsinov, V. Matveev, E. Miranda and S. Tabachnikov,
\emph{Open Problems, Questions, and Challenges in Finite-Dimensional
Integrable Systems}, Philos. Trans. R. Soc. A 376 (2018), 20170430.
\href{https://doi.org/10.1098/rsta.2017.0430}{doi:10.1098/rsta.2017.0430}.
Author versions: \href{https://arxiv.org/abs/1804.03737v2}{arXiv:1804.03737v2},
Problem 5.11, p.~29; v1, Problem 4.10, p.~37.
\bibitem{bkm} A. V. Bolsinov, A. Yu. Konyaev and V. S. Matveev,
\emph{Nijenhuis Geometry}, Adv. Math. 394 (2022), 108001.
\href{https://doi.org/10.1016/j.aim.2021.108001}{doi:10.1016/j.aim.2021.108001}.
Inspected author version: \href{https://arxiv.org/abs/1903.04603v2}{arXiv:1903.04603v2},
Theorems 3.1, 3.3, 4.5, 6.1 and p.~28 remark; publisher body not equated.
\bibitem{bkm3} A. V. Bolsinov, A. Yu. Konyaev and V. S. Matveev,
\emph{Nijenhuis geometry III: gl-regular Nijenhuis operators}, Rev. Mat.
Iberoam. 40 (2024), 155--188, online 3 March 2023.
\href{https://doi.org/10.4171/RMI/1416}{doi:10.4171/RMI/1416}, Theorem 1.10.
Author-v1: \href{https://arxiv.org/abs/2007.09506v1}{arXiv:2007.09506v1}, Theorem 1.4.
\bibitem{bkm4} A. V. Bolsinov, A. Yu. Konyaev and V. S. Matveev,
\emph{Nijenhuis geometry IV: conservation laws, symmetries and integration
of certain non-diagonalisable systems of hydrodynamic type in quadratures},
Nonlinearity 37 (2024), 105003.
\href{https://doi.org/10.1088/1361-6544/ad6acc}{doi:10.1088/1361-6544/ad6acc}.
Inspected author-v1: \href{https://arxiv.org/abs/2304.10626v1}{arXiv:2304.10626v1}.
\end{thebibliography}
\end{document}
